\input{./._relpath-to-latexroot.ltx}
\documentclass[\latexroot/\projectname]{subfiles}
\input{\latexroot/@resources/subfile-setup.ltx}
\begin{document}

\FloatBarrier
\hypertarget{robustness}{}

\section*{Robustness}
\whenintegrated{\label{sec:robustness}} % integrated doc: SST for labels

(This section was cut at the request of the editor, and is therefore not up to date with the latest versions of the code.  It is included here as a findable object in case anyone asks whether we have more robustness tests).

In this section, we analyze how sensitive our results are to some of the parameters in the model. 

\hypertarget{changing-the-interest-rate}{}\par\subsection*{Changing the interest rate}
\whenintegrated{\label{sec:robust_R}} % integrated doc: SST for labels 
In our baseline calibration, the interest rate is set to $1$ percent per quarter. 

\subsection*{Discount factor distributions with different interest rates}
\whenintegrated{\label{sec:robust_R_estim}} % integrated doc: SST for labels

Table~\ref{tab:robustness_R} shows the values we obtain for the discount factor distributions when we change the quarterly interest rate by either decreasing it to $0.5$ percent or increasing it to $1.5$ percent.

% Non-floating table (appendix is inside hiddencontent vbox environment)
\begin{center}
  \begin{tabular}{lc|cccccc}
    \toprule
                          &         & \multicolumn{2}{c}{Dropout} & \multicolumn{2}{c}{Highschool} & \multicolumn{2}{c}{College}                                    \\ \midrule
                          & Splurge & $\beta$                     & $\nabla$                       & $\beta$                     & $\nabla$    & $\beta$ & $\nabla$ \\ \midrule
    $R = 1.005$           & 0.307   & 0.740                       & 0.298                          & 0.927                       & $0.193^{*}$ & 0.989   & 0.0082   \\
    $R = 1.01$ (baseline) & 0.307   & 0.735                       & 0.298                          & 0.924                       & $0.137^{*}$ & 0.984   & 0.0096   \\
    $R = 1.015$           & 0.307   & 0.724                       & $0.357^{*}$                    & 0.919                       & $0.138^{*}$ & 0.979   & 0.0105
    \\ \bottomrule
  \end{tabular}
  \captionof{table}{Estimates of the splurge and $(\beta,\nabla)$ for each education group for different values of the interest rate $R$.
    A $*$ denotes an estimate of $\nabla$ that implies that the largest discount factor value in our discrete approximation would violate the GIC and is thus replaced with a value below the upper bound.}
  \label{tab:robustness_R}
\end{center}


\subsection*{Estimating discount factor distributions for different interest rates}
\whenintegrated{\label{app:DF_R}} % integrated doc: SST for labels

\subsection*{Results with different interest rates}
\whenintegrated{\label{sec:robust_R_results}} % integrated doc: SST for labels

In this section, we repeat the welfare analysis conducted in section~\ref{sec:welfare} for different values of the interest rate.

\FloatBarrier
\hypertarget{changing-risk-aversion}{}\par\subsection*{Changing risk aversion}
\whenintegrated{\label{sec:robust_gamma}} % integrated doc: SST for labels 

In our baseline calibration, consumers have a risk aversion of $\gamma=2$, which is quite common in macroeconomic models.

\subsection*{Discount factor distributions with different risk aversion}
\whenintegrated{\label{sec:robust_gamma_estim}} % integrated doc: SST for labels

Table~\ref{tab:robustness_gamma} displays the values we obtain for the splurge and the discount factor distributions when we change $\gamma$ to $1$ and $3$.

\subsection*{Results with different risk aversion}
\whenintegrated{\label{sec:robust_gamma_results}} % integrated doc: SST for labels

In this section, we conduct the welfare analysis for the case with $\gamma=3.0$.
\hypertarget{changing-benefits}{}\subsection*{Changing benefits}
\whenintegrated{\label{sec:robust_benefits}} % integrated doc: SST for labels 



\subsection*{Discount factor distributions with different benefits}
\whenintegrated{\label{sec:robust_benefits_estim}} % integrated doc: SST for labels

Table~\ref{tab:robustness_benefits} shows that when benefits are less generous and the precautionary motive is stronger, the estimated discount factor distributions are centered on lower values of~$\beta$.

\FloatBarrier
\subsection*{Results with different benefits}
\whenintegrated{\label{sec:robust_benefits_results}} % integrated doc: SST for labels

In this section, we perform the welfare analysis for different benefit replacement rates.

\FloatBarrier
\hypertarget{changing-the-properties-of-the-recession}{}\par\subsection*{Changing the properties of the recession}

In this section, we alter two properties of the recession and study the effect of those changes (one at a time) on our main results.


\FloatBarrier
%\hypertarget{changing-the-splurge}{}\par\subsection*{Changing the splurge}
%


\hypertarget{summary-of-robustness-exercises}{}\par\subsection*{Summary of robustness exercises}
\whenintegrated{\label{sec:robust_summary}} % integrated doc: SST for labels

The robustness checks in this section show that our main results are robust to a wide range of alternative parameter choices.


% Bibliography inclusion handled automatically by document structure
\smartbib

\end{document}
